\documentclass[11pt]{article}
\usepackage{amsmath,amssymb}
\usepackage[margin=2.5cm]{geometry}
\pagestyle{empty}

\begin{document}

\section*{Pair-level \texorpdfstring{$\Delta R$}{dR} trigger efficiency correction}

\subsection*{Single-muon efficiency}

The single-muon mu4 firing probability $\varepsilon(p_T, q\!\cdot\!\eta)$ is measured via tag-and-probe on muon pairs with $\Delta R > 0.8$, where proximity effects are negligible. The efficiency is extracted per centrality bin, fitted as a function of $p_T$ with a Fermi$+$log parametrization in each $q\!\cdot\!\eta$ slice.

\subsection*{Pair-level no-correlation efficiency}

Assuming the two muons fire independently (valid at large $\Delta R$), the probability that at least one muon in a pair fires mu4 is:
\begin{equation}
  \varepsilon_{\text{pair}}^{\text{nc}} = \varepsilon_1 + \varepsilon_2 - \varepsilon_1\,\varepsilon_2\,,
\end{equation}
where $\varepsilon_i \equiv \varepsilon(p_{T,i},\, q_i\!\cdot\!\eta_i)$ are the single-muon efficiencies evaluated at each muon's kinematics.

\subsection*{Inverse-weighted \texorpdfstring{$\Delta R$}{dR} ratio}

To isolate the $\Delta R$-dependent deviation from the no-correlation assumption, we form the ratio:
\begin{equation}
  R(\Delta R) = \frac{\displaystyle\sum_{\text{pairs},\;\Delta R\,\in\,\text{bin}}
    \frac{\mathbf{1}[\text{pair fires mu4}]}{\varepsilon_{\text{pair}}^{\text{nc}}}}
    {\displaystyle\sum_{\text{pairs},\;\Delta R\,\in\,\text{bin}} 1}\,.
\end{equation}
The denominator includes all pairs (no trigger requirement); the numerator counts pairs where at least one muon passes mu4, each weighted by $1/\varepsilon_{\text{pair}}^{\text{nc}}$.

If the no-correlation assumption held exactly and the inverse weighting were unbiased, $R(\Delta R) = 1$ everywhere. In practice, $R$ is flat at large $\Delta R$ but deviates from unity at small $\Delta R$ where trigger-object overlap suppresses the firing probability.

\subsection*{Why per-event inverse weighting}

An averaged correction---dividing by $\langle\varepsilon_{\text{pair}}^{\text{nc}}\rangle$ per $\Delta R$ bin---would leave the result entangled with the single-muon kinematic composition of each bin. Per-event inverse weighting avoids this: each pair's expected numerator contribution is
\begin{equation}
  E\!\left[\frac{\mathbf{1}[\text{fires}]}{\varepsilon_{\text{pair}}^{\text{nc}}}\right]
  = \frac{\varepsilon_{\text{pair}}^{\text{true}}}{\varepsilon_{\text{pair}}^{\text{nc}}} \,,
\end{equation}
which equals unity when the estimated efficiency matches the true one, \emph{regardless} of the kinematic distribution. The measurement is thus decorrelated from single-muon $p_T$ and $\eta$.

\subsection*{Large-$\Delta R$ offset and plateau normalization}

In practice, the flat-region value of $R$ is systematically $> 1$ ($\sim 1.2$). Since per-event weighting removes kinematic bias by construction, this offset reflects a systematic underestimate of $\varepsilon_{\text{pair}}^{\text{nc}}$ by the fitted parametrization (i.e.\ $\varepsilon^{\text{fit}} < \varepsilon^{\text{true}}$, so $\varepsilon^{\text{true}}/\varepsilon^{\text{fit}} > 1$). The bias is largest at low $p_T$ where the Fermi turn-on fit is most strained, consistent with the observed $p_T$- and centrality-dependence of the offset.

Since the physics of interest is the $\Delta R$ \emph{shape}---not the absolute scale---we extract the error-weighted plateau $\bar{R}$ from the flat region ($1.0 < \Delta R < 3.0$) and define the correction factor:
\begin{equation}
  \varepsilon_{\Delta R}^{\text{pair}}(\Delta R) = \frac{R(\Delta R)}{\bar{R}}\,.
\end{equation}

The full pair-level trigger efficiency applied to the data is then:
\begin{equation}
  P(\text{pair fires mu4}\mid\Delta R) = \varepsilon_{\text{pair}}^{\text{nc}}\;\times\;\varepsilon_{\Delta R}^{\text{pair}}(\Delta R)\,.
\end{equation}

\end{document}
